\documentclass[10pt,a4paper]{amsart}
\usepackage[margin=19mm]{geometry}
\usepackage{amsmath,amssymb,mathtools}
\usepackage[scr=rsfs,cal=euler]{mathalfa}
\usepackage{xfrac}
\usepackage[unicode,hidelinks,bookmarks=false]{hyperref}
\newcommand{\ie}{i.e.\ }
\newcommand{\cf}{cf.\ }
\newcommand{\resp}{respectively\ }
\newcommand{\Subsection}{\subsection}
\providecommand{\nobreakdash}{\nobreak\hbox{-}}
\setlength{\emergencystretch}{2em}
\multlinegap0pt
\usepackage{bm}
\usepackage{bbm}
\usepackage{dsfont}
\usepackage{xspace}
\usepackage{cancel}
\usepackage{mathtools}
\usepackage{amsthm}
\makeatletter
\@ifundefined{Theorem}{\newtheorem{Theorem}{Theorem}}{}
\@ifundefined{Proposition}{\newtheorem{Proposition}[Theorem]{Proposition}}{}
\makeatother
\title[Remarks on relative categoricity]{Remarks on relative categoricity}
\author{Anand Pillay}
\date{Source: arXiv:2602.05866v2 (CC BY 4.0)}
\begin{document}
\maketitle

\noindent\emph{Source and licence.} Remarks on relative categoricity. Version arXiv:2602.05866v2 (CC BY 4.0), distributed under the Creative Commons Attribution 4.0 International licence, as recorded in the arXiv licence metadata for this exact version and shown on the abstract page https://arxiv.org/abs/2602.05866v2. Full rights notice: licenses/arxiv-2602.05866-CC-BY-4.0.txt.\par\smallskip

\noindent\textit{Editorial scope.} Counted unit: the Proposition whose source label is \texttt{None} in the article (source0.tex lines 384--385, source label \texttt{None}), delivered together with the article's own complete original proof of it (source0.tex lines 386--421, one \texttt{proof} environment of 36 lines). No mathematics is written, repaired or re-derived: statement and proof are the article's own text line for line. Required context retained uncounted: none. Every definition, display and label of those lines is retained unchanged. Omitted from this package and not redistributed: 1--383, 422--555. Nothing inside a retained excerpt is omitted or altered: every retained body is delivered exactly as the article's lines are written, so no omission receipt of any kind accompanies this package. Citations: the retained excerpts carry no citation command at all, so no bibliography entry of the article is transcribed and no reference is redistributed. Reference adaptation: none was needed and none was made; every reference inside the delivered unit resolves inside it, so no unresolved reference or citation is produced. Printed slips kept literal: no printed word, symbol, hypothesis or number of the counted statement or of its proof was altered. Layout and class: the article's own class is \texttt{article} with packages including amsmath, amsfonts, amssymb, amsthm, amscd, xcolor; the wrapper is the lane's pinned \texttt{amsart} editorial wrapper at 10pt on A4, which restates the theorem-like environments the retained text needs and adds the article's own macro definitions that the retained text needs (none), with extra wrapper packages (bm, bbm, dsfont, xspace, cancel, mathtools). The delivered numbering is the unit's own. Assets: no retained span contains a figure environment, a graphics-inclusion command, a PDF-page inclusion, a table or a TikZ picture; no image file is copied into this package, the package declares no asset dependency and no image file is redistributed with it. Licence: version arXiv:2602.05866v2 of this article is distributed under the Creative Commons Attribution 4.0 International licence, as recorded in the arXiv licence metadata for this exact version and shown on the abstract page https://arxiv.org/abs/2602.05866v2. A full rights notice is delivered at licenses/arxiv-2602.05866-CC-BY-4.0.txt. Characters: every retained excerpt is pure ASCII, so the delivered engine, XeLaTeX with its default Latin Modern text font, reports no missing character.
\begin{Proposition} Suppose that $T$ is fully $\omega$-stable over $P$. Let $A$ be complete. Then the isolated types are dense in $S(A)$: for any consistent formula $\phi(x)$ over $A$, there is an isolated $p(x)\in S(A)$ containing $\phi(x)$.
\end{Proposition}

\begin{proof} We will assume the conclusion is false and use a routine argument to build continuum many good types over a countable complete subset of $A$, contradicting the $\omega$-stability assumption. 
So suppose $\phi(x)$ is a consistent formula over $A$ which is not in any isolated complete type over $A$.  We fix a bijection $f: \omega\setminus \{0\} \to \omega\times \omega$ such that if $f(n) = (r,i)$ then $r < n$. 

We build inductively formulas  $\psi_{\tau}(x)$ over $A$ for $\tau\in 2^{<\omega}$, countable complete subsets $A_{n}\subseteq A$ for $n<\omega$, and for each $n<\omega$ a list $\{\sigma_{(n,i)}(x): i<\omega\}$ of all formulas over $A_{n}$ of the form $\exists y \in P(\chi(x,y))$ with the properties::
\newline
(i)  $\psi_{<>}(x) = \phi(x)$, each $\psi_{\tau}(x)$ is consistent,   $\tau$ an initial segment of $\tau'$ implies $\models \psi_{\tau '}(x) \to \psi_{\tau}(x)$, and for each $\tau$ if $\tau_{1}$, $\tau_{2}$ are the $2$ successors of $\tau$ then $\psi_{\tau_{1}}(x)$ and $\psi_{\tau_{2}}(x)$ are contradictory,
\newline
(ii)  The $A_{n}$ are increasing and for each $n$ all the formulas $\psi_{\tau}(x)$ for $\tau$ of length $< n$  have parameters from $A_{n}$.
\newline
(iii)  Suppose $\ell(\tau) = n\geq 1$, and let $\sigma_{f(n)}(x)$ be $\exists y\in P (\chi(x,y))$. Then either $\models \psi_{\tau}(x) \to \neg\sigma_{f(n)}(x)$ or for some $c\in P(A_{n})$, $\models \psi_{\tau}(x) \to \chi(x,c)$. 

\vspace{2mm}
\noindent
To begin, let (by a downward Lowenheim-Skolem argument) $A_{0}$ be a countable complete subset of $A$ containing the parameters from $\phi(x)$. 
Suppose now $n< \omega$,  and $A_{r}$, $\psi_{\tau}(x)$ for $\ell(\tau) \leq n$ and also the list $\sigma_{r,i}(x)$ for $r\leq n$ and $i<\omega$, have been defined. 
Now for each $\tau$ of length $n$, $\psi_{\tau}(x)$ does not isolate a complete type over $A$ so there is a formula $\psi_{\tau}'(x)$ over $A$ such that both $\psi_{\tau}(x)\wedge \psi_{\tau}'(x)$ and $\psi_{\tau}(x)\wedge \neg\psi_{\tau}'(x)$ are consistent.   Let $\psi_{\tau ,0}''(x)  = \psi_{\tau}(x)\wedge \psi_{\tau}'(x)$ and $\psi_{\tau, 1}''(x) = \psi_{\tau}(x)\wedge \neg\psi_{\tau}'(x)$. 
So we have now defined $\psi_{\tau}''(x)$ for all $\tau$ of length $n+1$.  Now for $\tau$ of length $n+1$, if $\psi_{\tau}''(x)\wedge \sigma_{f(n+1)}(x)$ is consistent and $\sigma_{f(n+1}(x)$ is $(\exists y\in P)(\chi(x,y))$ then 
$(\exists y\in P)(\exists x)(\psi_{\tau}''(x) \wedge \chi(x,y))$ is over $A$ and consistent, so by completeness of $A$ pick $c_{\tau} \in P(A)$ such that $\psi_{\tau}''(x)\wedge \chi(x,c)$ is consistent.  Put $\psi_{\tau}(x) = \psi_{\tau}''(x)\wedge \chi(x,c)$.  Otherwise $\psi_{\tau}''(x)$ implies $\neg\sigma_{f(n+1)}(x)$, and just put $\psi_{\tau}(x) = \psi_{\tau}''(x)$. 

Now let $A_{n}'$ be $A_{n}$ together with the parameters in the formulas $\psi_{\tau}'(x)$ for $\tau$ of length $n$, and the $c_{\tau}$ for $\tau$ of length $n+1$. Let again by a  Lowenheim-Skolem argument, $A_{n+1}$ be a countable complete subset of $A$ containing $A_{n}'$.  Finally let  $\{\sigma_{n+1, i}(x): i<\omega\}$ be a list of all formulas over $A_{n+1}$ of the appropriate form $\exists y\in P(\chi(x,y))$.

\vspace{2mm}
\noindent
So we can carry out the inductive construction. Let $A'$ be the union of the $A_{n}$. So $A'$ is countable and complete. For each  $\eta\in 2^{\omega}$ let $p_{\eta}(x) = \{\psi_{\eta|n}(x): n <\omega\}$. So by (i), the $p_{\eta}$ are consistent sets of formulas over $A'$, and if $\eta \neq \eta'$, $p_{\eta}(x)$ and $p_{\eta'}(x)$ are contradictory.
 For each $\eta\in 2^{\omega}$ extend $p_{\eta}(x)$ to a complete type $p_{\eta}'(x)$ over $A'$.
\newline
{\em Claim.}  For $\eta\in 2^{\omega}$, $p_{\eta}'(x) \in S_{*}(A')$ namely is a good type over $A'$. 
\newline
{\em Proof of Claim.}  We must show that whenever a formula of the form $(\exists y\in P)(\chi(x,y))$ is in $p_{\eta}(x)$ then for some $c\in P(A')$, $\chi(x,c)\in p_{\eta}(x)$.  Now the formula $(\exists y\in P)(\chi(x,y))$ is over $A'$ so 
over $A_{n}$ for some $n$, so of the form $\sigma_{n,i}(x)$ for some $i$. Let $1\leq m < n$ be such that $f(m) = (n,i)$. Then $\psi_{\eta|m}(x) \in p_{\eta}'(x)$. So we cannot have that $\psi_{\eta|m}(x)$ implies $\neg\sigma_{n,i}(x)$.  So by (iii) 
of the construction $\models \psi_{\eta|m}(x) \to \chi(x,c)$ for some $c\in P(A_{m})\subset P(A')$. So $\chi(x,c) \in p_{\eta}'(x)$ and the Claim is proved.  \qed

So by the Claim and as the $p_{\eta}'(x)$ are distinct complete types over $A'$, we have that $S_{*}(A')$ is uncountable, contradicting the full $\omega$-stability over $P$ assumption.
This proves Proposition 2.8. 

\end{proof} 

\end{document}
